\chapter{Practical Control and Telemetry Examples}
\label{ch:practical-examples}

\section{Introduction}
\label{sec:practical-intro}

This chapter provides practical examples that demonstrate how the Android application
and the \textbf{MICSBOL/ESP32\_BT\_Controller} firmware interact in a real robot or embedded system.

The Android side continuously transmits a \textbf{control state} (sticks, knobs, and switches),
and also transmits \textbf{one-shot events} when the user presses buttons.

The ESP32 side:
\begin{itemize}
    \item decodes incoming packets over Bluetooth Serial (SPP),
    \item applies the received control values to real hardware outputs (PWM, direction pins, switches),
    \item and optionally sends telemetry back to the Android app (panels, indicators, plots).
\end{itemize}

\section{Control Data from the Android Application}
\label{sec:control-data}

\subsection{RC state packet (\texttt{rcStatePacket})}
The firmware stores the most recent control state in the structure defined in \texttt{packets.h}:

\begin{verbatim}
    rcStatePacket
\end{verbatim}

It contains the following fields:

\begin{itemize}
    \item \texttt{leftStickX}, \texttt{leftStickY}
    \item \texttt{rightStickX}, \texttt{rightStickY}
    \item \texttt{leftKnob}, \texttt{rightKnob}
    \item \texttt{switches} (bit mask)
\end{itemize}

In the repository code, the stick and knob channels are treated as 12-bit values and are mapped in the
range \textbf{0..4095} (see \texttt{control.cpp}, \texttt{hal\_output.cpp}).

\subsection{Event packet (\texttt{rcEventPacket})}
Buttons are transmitted as small event packets, and the event ID is passed to:

\begin{verbatim}
    controlHandleEvent(byte eventId)
\end{verbatim}

Events are useful for actions that should be \emph{momentary} (e.g., trigger, reset, horn) rather than a
continuous analog value.

\section{How the ESP32 Receives Packets}
\label{sec:receiver}

Incoming Bluetooth bytes are processed in:

\begin{verbatim}
    receiver.cpp  ->  handleBluetooth()
\end{verbatim}

The receiver maintains a rolling buffer and searches for packet headers.
In the current firmware implementation:

\begin{itemize}
    \item RC State header: \texttt{0xAA 0x55}
    \item Event header: \texttt{0xBB 0x66}
\end{itemize}

Once an event is decoded, \texttt{controlHandleEvent()} is called immediately.

\section{Where Control Values are Applied}
\label{sec:control-layer}

The main control loop is implemented in:

\begin{verbatim}
    control.cpp  ->  controlUpdate()
\end{verbatim}

This function reads values from \texttt{rcStatePacket} and sends them to the hardware abstraction layer
(\texttt{hal\_output.cpp}):

\begin{itemize}
    \item steering servo PWM
    \item left/right motor commands
    \item camera pan PWM
    \item LED brightness and buzzer intensity (PWM)
    \item switch outputs (bit mask decoded into per-channel outputs)
\end{itemize}

\section{Example 1: Steering Control Using the Left Joystick}
\label{sec:example-steering}

The firmware maps the left joystick horizontal axis to a servo-like PWM value:

\begin{verbatim}
    rcStatePacket.data.leftStickX
\end{verbatim}

In \texttt{control.cpp}, steering is generated by:

\begin{verbatim}
    static uint32_t mapServo(uint16_t v) {
        return map(v, 0, 4095, 205, 410);
    }

    uint32_t steerPWM = mapServo(rcStatePacket.data.leftStickX);
    halSetSteering(steerPWM);
\end{verbatim}

\subsection*{Practical interpretation}
\begin{itemize}
    \item The Android joystick produces a value from 0 to 4095.
    \item The ESP32 maps it into a smaller PWM range (205..410) suitable for its PWM configuration.
\end{itemize}

\noindent This is commonly used for:
\begin{itemize}
    \item RC car steering
    \item robot front-wheel steering
    \item a servo-actuated rudder
\end{itemize}

\section{Example 2: Differential Drive Motor Speed Control}
\label{sec:example-motors}

The firmware uses the vertical joystick axes as left/right motor channels:

\begin{verbatim}
    rcStatePacket.data.leftStickY
    rcStatePacket.data.rightStickY
\end{verbatim}

In \texttt{control.cpp}:

\begin{verbatim}
    uint32_t motorL = rcStatePacket.data.leftStickY;
    uint32_t motorR = rcStatePacket.data.rightStickY;

    halSetMotorLeft(motorL);
    halSetMotorRight(motorR);
\end{verbatim}

\subsection{Direction + speed mapping (implemented behavior)}
In \texttt{hal\_output.cpp}, each motor value is interpreted with a dead-band:

\begin{itemize}
    \item \textbf{0..2000:} direction pin LOW (reverse) and PWM mapped to 0..4095
    \item \textbf{2001..2094:} dead-band, motor target set to 0
    \item \textbf{2095..4095:} direction pin HIGH (forward) and PWM mapped to 0..4095
\end{itemize}

This design reduces jitter when the joystick is near center.

\subsection{Smooth ramping (advanced behavior)}
Instead of writing PWM instantly, the motor PWM ramps toward a target value. A timer interrupt updates
the motor outputs every 20 ms, moving speed by a fixed step (see \texttt{updateMotors()} in
\texttt{hal\_output.cpp}). This improves drivability and reduces current spikes.

\section{Example 3: Camera Pan Control}
\label{sec:example-camera}

The right joystick horizontal axis is used for camera pan:

\begin{verbatim}
    uint32_t panPWM = mapServo(rcStatePacket.data.rightStickX);
    halSetCameraPan(panPWM);
\end{verbatim}

\noindent Typical real-life uses:
\begin{itemize}
    \item inspection rover camera
    \item FPV camera positioning
    \item turret-like sensor mount
\end{itemize}

\section{Example 4: Using Knobs to Control LED and Buzzer PWM}
\label{sec:example-knobs}

The Android app provides two knobs, transmitted as 0..4095 values:

\begin{verbatim}
    rcStatePacket.data.leftKnob
    rcStatePacket.data.rightKnob
\end{verbatim}

In \texttt{control.cpp}, these map directly to PWM outputs via the HAL:

\begin{verbatim}
    halSetLed(rcStatePacket.data.leftKnob);
    halSetBuzzer(rcStatePacket.data.rightKnob);
\end{verbatim}

\subsection*{Practical uses}
\begin{itemize}
    \item LED brightness (headlights, underglow, work light)
    \item buzzer loudness / intensity (alerts, confirmation beeps)
    \item repurpose knobs for motor limit scaling or servo trim (custom firmware modification)
\end{itemize}

\section{Example 5: Controlling Digital Outputs with Switches}
\label{sec:example-switches}

Switches are transmitted as a single byte bit mask:

\begin{verbatim}
    byte sw = rcStatePacket.data.switches;
\end{verbatim}

In \texttt{control.cpp}, the firmware iterates through switch indices:

\begin{verbatim}
    for (int i = 0; i < 6; i++) {
        halSetSwitch(i, (sw >> i) & 1);
    }
\end{verbatim}

\subsection*{Important mode difference}
The number of switches supported depends on project mode (see \texttt{hal\_output.cpp}):

\begin{itemize}
    \item \textbf{FULL\_RC\_MODE:} only indices 0..3 are written to GPIO pins (4 switch outputs).
    \item \textbf{FULL\_RC\_MODE\_MCP:} indices 0..5 are written using MCP23017 outputs (6 switch outputs).
\end{itemize}

\noindent Real-world uses:
\begin{itemize}
    \item enable/arm switch (safety)
    \item headlights / beacon / underglow
    \item relay outputs (pump, solenoid, accessory power)
    \item mode selection (manual vs assisted)
\end{itemize}

\section{Example 6: Handling Button Events (Pulse Outputs)}
\label{sec:example-events}

Events are handled by:

\begin{verbatim}
    controlHandleEvent(byte eventId)
\end{verbatim}

In the repository implementation, events call \texttt{halPulseSwitch(index)}. The index mapping differs
by project mode. For example, in \texttt{FULL\_RC\_MODE}:

\begin{verbatim}
    case 0x01: halPulseSwitch(0); break;
    case 0x02: halPulseSwitch(1); break;
    case 0x03: halPulseSwitch(0); break;
    case 0x04: halPulseSwitch(1); break;
\end{verbatim}

\subsection{Non-blocking pulses}
\texttt{halPulseSwitch()} uses \texttt{pulseStart()} which generates a \textbf{non-blocking pulse}
(default 50 ms) and requires \texttt{pulseUpdate()} to be called in the main loop.

\noindent Typical event actions:
\begin{itemize}
    \item horn / short sound effect
    \item ``take picture'' trigger for a camera
    \item reset a subsystem (IMU re-zero, controller reset line)
    \item momentary relay (door latch, actuator kick)
\end{itemize}

\section{Telemetry Data Sent to the Android Application}
\label{sec:telemetry-data}

Telemetry is generated by:

\begin{verbatim}
    telemetry.cpp
    telemetry_source.cpp
\end{verbatim}

The design is intentionally split:
\begin{itemize}
    \item \textbf{\texttt{telemetry.cpp}} builds and sends packets (write-only).
    \item \textbf{\texttt{telemetry\_source.cpp}} provides values (hardware reads or debug injected values).
\end{itemize}

\subsection{Telemetry types and rates (implemented)}
\begin{itemize}
    \item \textbf{Indicators:} every 500 ms (\(\approx 2\,Hz\)).
    \item \textbf{Plots:} every 50 ms (\(\approx 20\,Hz\)).
    \item \textbf{Panels:} retransmitted briefly when values change (resend window for robustness).
\end{itemize}

\subsection{Config-on-connect}
On a new connection, the ESP32 sends a \textbf{configuration packet} that includes plot/panel/indicator
names (e.g., \texttt{Volts}, \texttt{Amps}, \texttt{RPMs}, \texttt{Left}, \texttt{Right}, \texttt{Throttle}, \texttt{Battery}).
This allows the Android app to label UI elements automatically.

\section{Example 7: Battery Telemetry as an Indicator}
\label{sec:example-battery}

The firmware’s indicator packet includes:
\begin{itemize}
    \item \texttt{analogValue} (0--100)
    \item \texttt{batteryLevel} (0--100)
    \item \texttt{digitalMask} (bitfield)
\end{itemize}

In normal mode, battery is read from an ADC pin and mapped to a percentage (see \texttt{hwReadIndicatorBattery()}
in \texttt{input\_hw.cpp}):

\begin{verbatim}
    return map(analogRead(PIN_ANALOG_IND2), 0, 4095, 0, 100);
\end{verbatim}

\noindent Practical wiring note: battery voltage must be scaled with a voltage divider so the ESP32 ADC
never exceeds 3.3V.

\section{Example 8: Telemetry Panels from Analog Inputs}
\label{sec:example-panels}

Panel telemetry uses two 16-bit values (\texttt{leftPanelValue}, \texttt{rightPanelValue}). The default
hardware implementation reads analog values and scales them from 0..4095 to 0..65535:

\begin{verbatim}
    uint16_t hwReadNumeric1() {
        uint32_t v = analogRead(PIN_NUMERIC1);
        return (uint16_t)((v * 65535UL) / 4095);
    }
\end{verbatim}

This is suitable for:
\begin{itemize}
    \item distance sensors with analog output
    \item potentiometer feedback
    \item scaled sensor signals (after conditioning)
\end{itemize}

\section{Example 9: Plot Signals (0..255 Fast Graphs)}
\label{sec:example-plots}

Plots transmit three 8-bit values (0..255) at \(\approx 20\,Hz\). The default hardware implementation reads
ADC values and maps them to 0..255:

\begin{verbatim}
    uint8_t hwReadPlot1() { return map(analogRead(PIN_NUMERIC1), 0, 4095, 0, 255); }
    uint8_t hwReadPlot2() { return map(analogRead(PIN_NUMERIC2), 0, 4095, 0, 255); }
    uint8_t hwReadPlot3() { return map(analogRead(PIN_ANALOG_IND2), 0, 4095, 0, 255); }
\end{verbatim}

\noindent Practical plot ideas:
\begin{itemize}
    \item battery voltage trend (sag under acceleration)
    \item current sensor trend (detect stalls)
    \item RPM/encoder-derived speed trend
\end{itemize}

\section{Example 10: Digital Mask Indicators}
\label{sec:example-digital-mask}

The indicator telemetry includes a digital mask byte that can represent up to 8 on/off signals.
In \texttt{input\_hw.cpp}, a mask is built from indicator GPIO reads (or from MCP23017 inputs in MCP mode).

This is useful for compactly showing:
\begin{itemize}
    \item limit switches
    \item fault flags
    \item bumper sensors
    \item ``armed'' / ``ready'' signals
\end{itemize}

\section{Summary}
\label{sec:practical-summary}

The system supports bidirectional communication:

\begin{itemize}
    \item \textbf{Android \(\rightarrow\) ESP32:} continuous RC state + discrete events.
    \item \textbf{ESP32 \(\rightarrow\) Android:} panels, indicators, and plots with automatic label configuration.
\end{itemize}

By combining control inputs (sticks, knobs, switches, events) with live telemetry feedback, you can build
robots and embedded systems that are both \textbf{easy to operate} and \textbf{easy to debug in real time}.